\documentclass[12pt]{article}

\input{../../_assets/preambulo.tex}


\begin{document}

    % 1. Foto de fondo
    % 2. Título
    % 3. Encabezado Izquierdo
    % 4. Color de fondo
    % 5. Coord x del titulo
    % 6. Coord y del titulo
    % 7. Fecha

    
    \input{../../_assets/portada}
    \portadaExamen{ffccA4.jpg}{Ecuaciones\\Diferenciales II\\Examen XII}{Ecuaciones Diferenciales II. Examen XII}{MidnightBlue}{-8}{28}{2026}{}

    \begin{description}
        \item[Asignatura] Ecuaciones Diferenciales II.
        \item[Curso Académico] 2014/2015.
        \item[Grado] Doble Grado en Ingeniería Informática y Matemáticas.
        \item[Grupo] Único.
        \item[Profesor] Rafael Ortega Ríos.
        \item[Descripción] Convocatoria Ordinaria.
        \item[Fecha] 26 de junio de 2015.
        % \item[Duración] ---.
    
    \end{description}
    \newpage


    % ------------------------------------
    
    \begin{ejercicio}
        Se considera el sistema
        \[ x(t) = 7 + \int_{0}^{t} y(s) \, ds, \quad y(t) = -\int_{0}^{t} x(s) \, ds. \]
        ¿Existe solución? ¿Es única? En caso afirmativo calcula dicha solución.
    \end{ejercicio}

    \begin{ejercicio}
        Encuentra dos soluciones maximales distintas del problema
        \[ \dot{x} = +\sqrt{|x|}, \quad x(0) = 1. \]
        Demuestra que hay unicidad local.
    \end{ejercicio}

    \begin{ejercicio}
        Se considera la sucesión de funciones $f_n: \bb{R} \to \bb{R}$,
        \[ f_n(t) = \begin{cases} -1 & \text{si } t \leq -n \\ \nicefrac{t}{n} & \text{si } |t| < n \\ 1 & \text{si } t \geq n. \end{cases} \]
        Demuestra que la sucesión es equicontinua. ¿Es uniformemente acotada? Calcula $\lim\limits_{n \to \infty} f_n(t) = f(t)$. ¿Hay convergencia uniforme de $f_n$ a $f$?
    \end{ejercicio}

    \begin{ejercicio}
        Demuestra que todas las soluciones maximales de la ecuación
        \[ \ddot{x} + x^5 = 0 \]
        están definidas en $\left]-\infty, +\infty\right[$.
    \end{ejercicio}

    \begin{ejercicio}
        Dado $h \in \bb{R}$, la solución de
        \[ \dot{x} = \operatorname{sen} x, \quad x(0) = \pi + h \]
        se denota por $x(t; h)$. Calcula
        $$\Delta(t) = \lim\limits_{h \to 0} \frac{1}{h} \left[ x(t; h) - x(t; 0) \right].$$
    \end{ejercicio}

    \begin{ejercicio}
        Se supone que $f: \left[0, \infty\right[ \to \bb{R}$ es una función continua que cumple
        \[ f(t) \leq 5e^{-t} + 3 \int_{0}^{t} e^{s-t} f(s) \, ds \]
        si $t \geq 0$. Encuentra números $a$ y $b$ para los que, si $t\geq 0$, se cumpla
        \[ f(t) \leq a e^{bt}. \]
    \end{ejercicio}

    \begin{ejercicio}
        En el espacio de las matrices $\bb{R}^{4 \times 4}$ se considera la norma dada por $\|A\| = \max\limits_{1 \leq i \leq 4} \sum\limits_{j=1}^{4} |a_{ij}|$. Encuentra un polinomio $p(t)$ para el que se cumpla la desigualdad
        \[ \|e^{Bt}\| \leq p(t) e^{3t} \]
        si $t \geq 0$, donde $B$ es la matriz
        \[ B = \begin{pmatrix} 3 & 1 & 0 & 0 \\ 0 & 3 & 1 & 0 \\ 0 & 0 & 3 & 1 \\ 0 & 0 & 0 & 3 \end{pmatrix}. \]
    \end{ejercicio}

    \begin{ejercicio}
        Estudia la estabilidad del equilibrio $x = y = 0$ para el sistema
        \[ \dot{x} = y^5, \quad \dot{y} = -x^7. \]
    \end{ejercicio}

    % \newpage
    % \setcounter{ejercicio}{0} % Reiniciar contador de ejercicios
    % \noindent
    % \textbf{Solución.}

\end{document}